\DocumentMetadata{tagging = off}

\documentclass{ltx-talk}




% Use STIX Two Math to match the "Times" look for your equations








\EditInstance{footer}{std}{
  color         = black,
  separator     = \hfill|\hfill,
  element-order = {framenumber},
  %element-order    = {date, author, title, institute, framenumber},
  right-hspace  = 3mm,
  left-hspace   = 3mm,
}


%\usepackage{amsmath}
%\newtheorem{definition}{Definition}
\definecolor{myblue}{rgb}{0.0, 0.1, 0.8}


% Preamble
\newtheorem{definition}{\color{myblue}Definition}
\newtheorem{theorem}{\color{myblue}Theorem}
\newtheorem{corollary}{\color{myblue}Corollary}
\newtheorem{remark}{\color{myblue}Remark}
\newtheorem{example}{\color{myblue}Example}



\usepackage[ruled,vlined]{algorithm2e}
\SetNlSty{textbf}{\color{red}}{}








\usepackage[ruled,vlined]{algorithm2e}
\SetNlSty{textbf}{\color{red}}{} % Colors line numbers to match the class
% 3. Define a simple command for Chinese text

%\newfontfamily\zhfont{FandolSong-Regular.otf}[
 % Script=CJK,
 % Renderer=HarfBuzz % Use the faster HarfBuzz engine in LuaLaTeX
%]
%\newcommand{\zh}[1]{{\zhfont #1}}


%\setmainfont{Times New Roman}

%\setmathfont{STIX Two Math}


%\usepackage[fontset=fandol]{ctex}

\pagecolor{green!5}



\title{the first presentation of Hazhy}
\subtitle{the first presentation of Hazhy}
\author{Hazhy}
\institute{None}
\date{2026-1-13}



\begin{document}
\maketitle

\begin{frame}
\frametitle{Sparsity and Manifolds}
    % English and Math stay in the default New Computer Modern Sans/Math
    The sparsity is defined by the $\ell_1$-norm:
    \[ \min \|x\|_{\ell_1} \quad \text{subject to} \quad Ax = b \]

    \begin{itemize}
        \item \textbf{Riemannian Manifold:}
        \item \textbf{Optimization:}
        \item  
    \end{itemize}

    % Test the \ell symbol:
    The cursive $\ell$ is preserved from the default NewCMSansMath.
\end{frame}
\begin{frame}

\begin{definition}
        A vector $x \in \mathbb{R}^n$ is $s$-sparse if it has at most $s$ non-zero entries.
    \end{definition}
  \begin{theorem}
          A vector $x \in \mathbb{R}^n$ is $s$-sparse if it has at most $s$ non-zero entries.
  \end{theorem}
    \begin{description}
        \item[as]  a test
    \end{description}
     \begin{corollary}
         A vector $x \in \mathbb{R}^n$ is $s$-sparse if it has at most $s$ non-zero entries.
     \end{corollary}
     \begin{remark}
         A vector $x \in \mathbb{R}^n$ is $s$-sparse if it has at most $s$ non-zero entries.
     \end{remark}
     \begin{example}
      
     \end{example}
\end{frame}


\begin{frame*}{Sample Code}

\begin{verbatim}
print("Hello, sparse recovery!")
\end{verbatim}


\end{frame*}




\begin{frame}
\frametitle{jjj}

\end{frame}

\begin{frame}{Algorithm: Basis Pursuit}
\begin{algorithm}[H]
    \caption{L1-Minimization}
    \KwIn{Matrix $A$, vector $b$}
    \KwOut{Sparse vector $x$}
    Solve $\min \|x\|_1$ subject to $Ax = b$\;
\end{algorithm}
\end{frame}

\end{document}
